\section{Probability}

\answer{PG1-1}
Let $E$ denote at least one head and $C$ the coin type. Conditional on type $i$, the complement of $E$ is $n$ consecutive tails, so its likelihood is $L_i=1-(1-p_i)^n$. Bayes' formula gives
\[
\boxed{w:=\Prob(C=1\mid E)=\frac{\pi L_1}{\pi L_1+(1-\pi)L_2}.}
\]
The next toss is independent of $E$ conditional on the type. Thus
\[
\boxed{\Prob(H_{n+1}\mid E)=wp_1+(1-w)p_2.}
\]
For $n>1$, the record combines many sequences and has likelihood $L_i$, generally not the single-head likelihood $p_i$. As $n\to\infty$, both likelihoods approach one and the posterior returns to the prior: either coin eventually produces a head with probability approaching one.

\answer{PG1-2}
Put $s=\alpha+\beta$ and let $E$ be the conditioning event. The joint density that $A$ rings at $u$ while $B$ has not rung is $\alpha e^{-su}$. Hence
\[
\Prob(E)=\int_0^t\alpha e^{-su}\dd u=\frac\alpha s(1-e^{-st}),
\]
and division gives
\[
\boxed{f_{T\mid E}(u)=\frac{s e^{-su}}{1-e^{-st}},\quad0<u<t.}
\]
Integration by parts yields
\[
\E(T\mid E)=\frac{-t e^{-st}+(1-e^{-st})/s}{1-e^{-st}}
=\boxed{\frac1s-\frac{t}{e^{st}-1}}.
\]
As $t\to\infty$, the density becomes exponential with rate $s$ and the mean tends to $1/s$. Knowing the winning clock alone does not change the first ringing time; the additional deadline truncates it.

\answer{PG1-3}
Put $d=1-a-b$. After removing the common multinomial coefficient, the likelihoods under $T=1$ and $T=0$ are $a^hb^zd^e$ and $b^ha^zd^e$. Therefore
\[
\boxed{O:=\frac{\Prob(T=1\mid\text{record})}{\Prob(T=0\mid\text{record})}
=\frac\pi{1-\pi}\left(\frac ab\right)^{h-z}.}
\]
The posterior probability is $O/(1+O)$. The factor $d^e$ cancels, so for fixed $h,z$ the number of missing reports is irrelevant in this model.

If missing probabilities under the two truths are $d_1,d_0$, each missing report instead contributes likelihood ratio $d_1/d_0$. The probabilities of nonmissing outputs must also be recomputed under the revised model. Missingness can therefore be informative and cannot simply be discarded.

\answer{PG1-4}
Let $S=(X+Z)/\sqrt2$, $D=(X-Z)/\sqrt2$, and $b=a/\sqrt2$. The variables $S,D$ are independent standard normals, and selection restricts only $|S|<b$, an event of probability $2\Phi(b)-1$. The conditional mean is zero by symmetry. Integration by parts gives
\[
v:=\Var(S\mid |S|<b)
=\frac{\int_{-b}^b s^2\phi(s)\dd s}{2\Phi(b)-1}
=1-\frac{2b\phi(b)}{2\Phi(b)-1}.
\]
The variance of $D$ remains one, and it remains independent of the truncated $S$. Thus
\[
\boxed{\operatorname{Corr}(X,Z\mid |X+Z|<a)=\frac{v-1}{v+1}.}
\]
As $b\downarrow0$, the bound $|S|<b$ implies $0\le v\le b^2$, so the correlation tends to $-1$. As $b\to\infty$, truncation disappears, $v\to1$, and the correlation tends to zero.

\answer{PG2-1}
The first draw cannot repeat, and the fourth must repeat, so $T\in\{2,3,4\}$. The probability that the first two types differ is
\[
\Prob(T>2)=1-\left[(1/2)^2+2(1/4)^2\right]=\frac58.
\]
Three distinct draws require one of each type in one of $3!$ orders, giving
\[
\Prob(T>3)=6\cdot\frac12\cdot\frac14\cdot\frac14=\frac3{16}.
\]
Since $\Prob(T>0)=\Prob(T>1)=1$, the tail sum yields
\[
\boxed{\E T=1+1+\frac58+\frac3{16}=\frac{45}{16},
\qquad \Prob(T=4)=\frac3{16}.}
\]
The last equality uses $T\le4$; it is not the probability of any specified type appearing.

\answer{PG2-2}
Each red run has exactly one black-to-red boundary. Any directed circular adjacency has a black start and red end with probability $nm/[N(N-1)]$. Summing over all $N$ boundaries gives
\[
\boxed{\E(\text{red runs})=\frac{mn}{N-1}.}
\]
If all red cards are consecutive, specifying the start of the unique red run determines the entire color arrangement. There are $N$ choices, producing distinct labeled arrangements because both colors occur. Since there are $\binom Nm$ arrangements in total,
\[
\boxed{\Prob(\text{one red run})=\frac N{\binom Nm}.}
\]
For $m=1$ or $n=1$, this probability is one, as it should be.

\answer{PG2-3}
For each triple $A$, let $I_A$ indicate that its objects share a bucket. Then $\E I_A=s_3$, so $\E C=\binom m3s_3$. Two disjoint triples are independent. Triples sharing two objects collide together exactly when all four objects share a bucket, giving covariance $s_4-s_3^2$. If they share one object, all five must share a bucket, giving covariance $s_5-s_3^2$.

There are $6\binom m4$ unordered pairs sharing two objects: choose their common pair among four objects. There are $15\binom m5$ pairs sharing one object: choose that object in five ways, then partition the remaining four into two unordered pairs in three ways. Each unordered pair contributes twice to the variance expansion. Thus
\[
\boxed{\begin{aligned}
\E C&=\binom m3s_3,\\
\Var C&=\binom m3s_3(1-s_3)
+12\binom m4(s_4-s_3^2)\\
&\quad+30\binom m5(s_5-s_3^2).
\end{aligned}}
\]
With $\binom mr=0$ for $r>m$, the formula includes $m=3,4$.

\answer{PG2-4}
Let $Y_i$ equal $1$ if ordinary card $i$ precedes every target, $-1$ if it follows every target, and zero otherwise. Reversing the permutation gives $\E Y_i=0$. Among this card and the $K$ targets, each extreme position has probability $1/(K+1)$, so $\E Y_i^2=2/(K+1)$.

For $i\ne j$, the probability that both precede all targets is $2/[(K+1)(K+2)]$, as is the probability that both follow them. The probability that specified $i$ is before and $j$ after is $1/[(K+1)(K+2)]$, and likewise in reverse. Hence
\[
\E(Y_iY_j)=\frac2{(K+1)(K+2)}.
\]
Expanding $D=\sum_iw_iY_i$ and using $2\sum_{i<j}w_iw_j=(\sum_iw_i)^2-\sum_iw_i^2$ gives
\[
\boxed{\E D=0,\qquad
\Var D=\frac{2[(K+1)\sum_iw_i^2+(\sum_iw_i)^2]}{(K+1)(K+2)}.}
\]
The algebra does not require positive weights. For $m=0$, all sums and the variance are zero.

\answer{PG3-1}
The simplex volume gives the waiting-time tail:
\[
\Prob(N>k)=\Prob(U_1+\cdots+U_k\le1)=\frac1{k!},\quad k\ge0.
\]
Thus $\E N=\sum_{k\ge0}1/k!=e$. Writing the cost as a nonnegative series permits interchange of expectation and summation:
\[
\E C=\sum_{j\ge1}\E[U_j^2\ind\{N\ge j\}].
\]
The event $\{N\ge j\}$ says that the first $j-1$ increments sum to at most one. It depends only on the past and is independent of the fresh $U_j$. Each term is therefore $\E U_j^2\Prob(N\ge j)=\Prob(N\ge j)/3$, so
\[
\boxed{\E C=\frac13\E N=\frac e3.}
\]
The stopping time does depend on the increment sequence. Only the decision whether draw $j$ occurs is independent of that draw's as-yet unseen increment.

\answer{PG3-2}
Put $t=-\log x$ and $b=\log2$. The negative log-length is $nb+G_n$, where $G_n=\sum_{i=1}^n(-\log U_i)$ has density $s^{n-1}e^{-s}/(n-1)!$. Since $G_n>0$ almost surely, crossing has occurred whenever $nb\ge t$. Thus $N\le L:=\lceil t/b\rceil$ almost surely.

For $n\ge1$, set $a_n=t-nb$. If $a_n\le0$, the tail is zero; otherwise,
\[
\boxed{\Prob(N>n)=\Prob(G_n\le a_n)
=1-e^{-a_n}\sum_{j=0}^{n-1}\frac{a_n^j}{j!}.}
\]
The final expression follows by repeated integration by parts of the gamma density. Including $\Prob(N>0)=1$ gives
\[
\boxed{\E N=1+\sum_{n=1}^{L-1}
\left[1-e^{-(t-nb)}\sum_{j=0}^{n-1}\frac{(t-nb)^j}{j!}\right].}
\]
For $x=1/4$, $L=2$. Crossing on the first draw means $U_1/2<1/4$, which has probability $1/2$; otherwise crossing occurs on the second. Hence $\Prob(N=1)=\Prob(N=2)=1/2$ and $\E N=3/2$.

\answer{PG3-3}
For $0<s<1$, summing the densities of all positive partial sums gives $u(s)=\sum_{k\ge1}s^{k-1}/(k-1)!=e^s$. A first entry into $(a,b]$ at $s$ may follow $S_0=0$ or a positive partial sum $t\le a$. Since $s-t\in(0,1)$, the next-increment density is one. The first-entry density is therefore
\[
1+\int_0^a e^t\dd t=e^a,\quad a<s\le b.
\]
Put $d=b-a$. Exactly $k\ge1$ entries have ordered positions $a<s_1<\cdots<s_k\le b$. Internal transition densities are one, and the probability of jumping past $b$ after $s_k$ is $1-b+s_k$. For $s_k=a+v$, the earlier $k-1$ positions have simplex volume $v^{k-1}/(k-1)!$. Consequently,
\[
\begin{aligned}
\Prob(K=k)
&=\frac{e^a}{(k-1)!}\int_0^d v^{k-1}(1-d+v)\dd v\\
&=\boxed{e^a\left(\frac{d^k}{k!}-\frac{d^{k+1}}{(k+1)!}\right)}.
\end{aligned}
\]
Summing over $k\ge1$ telescopes to $e^ad$, so $\boxed{\Prob(K=0)=1-e^ad}$. This is nonnegative because $d\le1-a$ and $e^a(1-a)\le1$. Finally, integrate the renewal density to obtain
\[
\boxed{\E K=\int_a^bu(s)\dd s=e^b-e^a.}
\]

\answer{PG3-4}
Future shrinkage factors are independent of the observed history $\mathcal F_n$. The normalized representation in Problem \ref{prob:P24} gives
\[
Q=Q_n+(1-R_n)D',
\]
where $D'$ is independent of $\mathcal F_n$, with mean one and variance $1/2$. Hence
\[
\boxed{\widehat Q=\E(Q\mid\mathcal F_n)=Q_n+1-R_n,}
\quad
\boxed{\Var(Q\mid\mathcal F_n)=\frac{(1-R_n)^2}{2}.}
\]
For any square-integrable predictor $Z$ based on the history,
\[
\E(Q-Z)^2=\E\Var(Q\mid\mathcal F_n)+\E(\widehat Q-Z)^2.
\]
The cross term vanishes by conditional expectation, proving optimality of $\widehat Q$. Since $1-R_n=(1-x)\prod_{i=1}^nU_i$,
\[
\boxed{\E(Q-\widehat Q)^2=\frac{(1-x)^2}{2\cdot3^n}.}
\]
The partial sum omits a positive tail with conditional mean $1-R_n$. Adding that mean removes the predictable bias.

\answer{PG4-1}
Introduce $S=U+V$. Then $U=WS$, $V=(1-W)S$, with Jacobian $s$. The unit-square constraints become
\[
0<w<1,\qquad 0<s<\frac1{\max(w,1-w)}.
\]
Integrating out $s$ gives
\[
\boxed{f_W(w)=\int_0^{1/\max(w,1-w)}s\dd s
=\frac1{2\max(w,1-w)^2},\quad0<w<1.}
\]
The density is not constant, so the distribution is not uniform. Exchanging $U,V$ replaces $W$ by $1-W$, implying $\E W=1/2$. For the second moment,
\[
\begin{aligned}
\E W^2&=\frac12\int_0^{1/2}\frac{w^2}{(1-w)^2}\dd w
+\frac12\int_{1/2}^1 1\dd w\\
&=\left(\frac34-\log2\right)+\frac14=1-\log2.
\end{aligned}
\]
Thus $\boxed{\Var W=3/4-\log2}$. Symmetry of the two shares determines the mean, not a constant density.

\answer{PG4-2}
For $0<l<h<1$, the joint density of the extremes is $n(n-1)(h-l)^{n-2}$: select the minimum and maximum observations, then place the other $n-2$ between them. Set $r=h-l$ and $c=(l+h)/2$. The inverse Jacobian is one and $0<r<2\min(c,1-c)$. Hence
\[
\boxed{f_C(c)=\int_0^{2\min(c,1-c)}n(n-1)r^{n-2}\dd r
=n[2\min(c,1-c)]^{n-1},\quad0<c<1.}
\]
Symmetry about $1/2$ gives $\E C=1/2$. Using that symmetry and substituting $u=2c$,
\[
\begin{aligned}
\Var C&=2\int_0^{1/2}(1/2-c)^2n(2c)^{n-1}\dd c\\
&=\frac n4\int_0^1(1-u)^2u^{n-1}\dd u
=\boxed{\frac1{2(n+1)(n+2)}}.
\end{aligned}
\]
The last integral follows by expanding $(1-u)^2$, giving $1/n-2/(n+1)+1/(n+2)$.

\answer{PG4-3}
Let the spacings be $D_1,\ldots,D_n$. Given them, the random point falls in spacing $i$ with probability $D_i$, not $1/n$. For any nonnegative function $g$,
\[
\E g(L)=\E\sum_iD_i g(D_i)=n\E[D_1g(D_1)].
\]
One spacing has tail $\Prob(D_1>l)=(1-l)^{n-1}$ and density $(n-1)(1-l)^{n-2}$. Therefore
\[
\boxed{f_L(l)=n(n-1)l(1-l)^{n-2},\quad0<l<1.}
\]
This is the $\operatorname{Beta}(2,n-1)$ density. Beta integration gives
\[
\boxed{\E L=\frac2{n+1},\qquad
\E(1/L)=n(n-1)\int_0^1(1-l)^{n-2}\dd l=n.}
\]
Choosing a spacing uniformly by index gives mean $\E(\sum_iD_i/n)=1/n$, strictly smaller than $2/(n+1)$. A random point is more likely to select a longer spacing.

\answer{PG4-4}
The vector's direction is uniform on the circle. Projections along two directions separated by $\theta$ have opposite signs in two sectors of total angle $2\theta$. Thus $\Prob(D_1=1)=\Prob(D_2=1)=\theta/\pi$.

Let $V_0,V_1,V_2$ be the actual projection values. A trigonometric identity gives
\[
V_1=\frac{V_0+V_2}{2\cos\theta}.
\]
Since $\cos\theta>0$, equal signs of the two outer projections force the middle sign to agree. Two successive sign changes are therefore impossible: $\Prob(D_1=D_2=1)=0$. Zero projections have probability zero. Combining with the marginals,
\[
\boxed{\begin{array}{c|cc}
& D_2=0&D_2=1\\ \hline
D_1=0&1-2\theta/\pi&\theta/\pi\\
D_1=1&\theta/\pi&0
\end{array}}
\]
Consequently,
\[
\boxed{\Cov(D_1,D_2)=-\frac{\theta^2}{\pi^2},\qquad
\Prob(A=B=C)=1-\frac{2\theta}{\pi}.}
\]
The adjacent comparisons share a middle projection, so their sign changes are not independent.

\par\Needspace{8\baselineskip}
\subsection*{Section Exercises}
\answer{OP02}
Write $Z$ for B's head count in the first $2n$ tosses and $B$ for the indicator of a head on the final toss. Then $Y=Z+B$, with $B$ independent and equally likely to be $0$ or $1$. Put $D=Z-2X$. Complement every outcome in A's tosses and B's first $2n$ tosses. This preserves probability and maps
\[
D\longmapsto(2n-Z)-2(n-X)=-D.
\]
Thus $\Prob(D>0)=\Prob(D<0)$. A positive $D$ guarantees success; a negative $D$ cannot be made strictly positive by adding $B\le1$; and $D=0$ gives success exactly when $B=1$. Consequently,
\[
\Prob(Y>2X)=\Prob(D>0)+\frac12\Prob(D=0)
=\boxed{\frac12}.
\]
The symmetry is that of the difference $Z-2X$; it does not assert that $Z$ and $2X$ have the same distribution.

\answer{OP30}
Let $F$ mean that the double-headed coin was selected, and let $H$ denote two heads. The prior probabilities and likelihoods are
\[
\Prob(F)=\Prob(F^c)=\frac12,\qquad
\Prob(H\mid F)=1,\quad \Prob(H\mid F^c)=\frac14.
\]
The law of total probability gives $\Prob(H)=\frac12+\frac12\cdot\frac14=\frac58$. Bayes' formula therefore yields
\[
\Prob(F\mid H)=\frac{\Prob(F)\Prob(H\mid F)}{\Prob(H)}
=\frac{1/2}{5/8}=\boxed{\frac45}.
\]
Both observations come from the same selected coin, so the coin type must not be resampled between tosses.

\answer{OP18}
Regard each connected component that has not yet closed as an open chain. Initially there are $n$ such chains, each with two free ends. With $k$ open chains there are $2k$ free ends. Joining the two ends of one chain closes a loop; joining ends of different chains merges them. Either operation reduces the number of open chains by one. The process therefore visits $k=n,n-1,\ldots,1$ exactly once each.

Of the $\binom{2k}{2}$ possible endpoint pairs, exactly $k$ join the two ends of the same chain. Thus the conditional probability of creating a loop at that step is always
\[
\frac{k}{\binom{2k}{2}}=\frac1{2k-1},
\]
regardless of the previous chain configuration. If $I_k$ indicates a loop at that step, the final loop count is $C=\sum_{k=1}^nI_k$. Linearity of expectation gives
\[
\boxed{\E C=\sum_{k=1}^n\frac1{2k-1}.}
\]
Independence of the loop indicators is not needed for this expectation.

\answer{OP25}
The four endpoints of any two chords are distinct almost surely. Conditional on their cyclic order, the three unordered pairings into two chords are equally likely. Only the pairing that joins alternating endpoints produces an interior crossing. Hence each pair of chords intersects with probability $1/3$.

Let $I_{ij}$ indicate an interior crossing of chords $i$ and $j$. Coincident endpoints, coincident chords, and triple intersections have probability zero under continuous sampling. For example, after fixing an intersection of two chords and one endpoint of a third chord, at most one position for its other endpoint makes it pass through that intersection. Thus the number of distinct crossings is almost surely
\[
K=\sum_{1\le i<j\le n}I_{ij}.
\]
Linearity of expectation, without any independence assumption on the indicators, gives
\[
\boxed{\E K=\binom n2\frac13=\frac{n(n-1)}6.}
\]
The sampling rule is essential: a random chord here is defined by two independent uniform endpoints.

\answer{OP26}
Count the non-aces left after the last ace. For each non-ace $i$, let $I_i$ indicate that it appears after all four aces. The relative order of that card and the four aces is uniform, so the specified non-ace is last among those five cards with probability $1/5$. Hence $\E I_i=1/5$.

There are $48$ non-aces. If $R$ is the number of cards after the last ace, then
\[
R=\sum_{i=1}^{48}I_i,\qquad \E R=\frac{48}{5}.
\]
Since $T+R=52$, the answer is
\[
\boxed{\E T=52-\frac{48}{5}=\frac{212}{5}=42.4.}
\]
This is the stopping position including the fourth ace. The number of cards strictly before that ace would be smaller by one.

\answer{OP22}
While the pattern is still valid and the current value is $x$, let $f(x)$ be the expected number of additional draws when the next comparison must be an increase; let $g(x)$ be the corresponding expectation when a decrease is required. One draw is always made, and continuation occurs only if that comparison succeeds. Thus
\[
f(x)=1+\int_x^1g(y)\dd y,\qquad
 g(x)=1+\int_0^xf(y)\dd y.
\]
The expectations are finite. If the first $n$ values obey the pattern, every disjoint pair $(U_1,U_2),(U_3,U_4),\ldots$ must increase, so the survival probability is at most $2^{-\lfloor n/2\rfloor}$. The same geometric bound controls the remaining draws from any state and justifies the integral equations.

Differentiating gives $f'=-g$ and $g'=f$, hence $f''+f=0$. The boundary conditions $f(1)=1$ and $g(0)=1$, or $f'(0)=-1$, yield
\[
f(x)=A\cos x-\sin x,\qquad
A=\frac{1+\sin1}{\cos1}=\sec1+\tan1.
\]
After the first draw, the next comparison must be an increase. Therefore,
\[
\begin{aligned}
\E T&=1+\int_0^1f(x)\dd x\\
&=A\sin1+\cos1
=\boxed{\sec1+\tan1}.
\end{aligned}
\]
The trigonometric arguments are in radians.

\answer{OP24}
Start more generally from $R_0=x\in[0,1)$. The process can be represented as
\[
1-R_k=(1-x)\prod_{j=1}^kU_j,
\qquad U_j\overset{\mathrm{iid}}\sim\operatorname{Unif}(0,1).
\]
The finite products $Q_n=\prod_{k=1}^nR_k$ decrease and lie in $[0,1]$, so they converge to a limit $Q$. Set $g(x)=\E[Q\mid R_0=x]$ and define $g(1)=1$. Bounded convergence makes the requested limit equal to $g(0)$.

The right boundary condition needs justification. The inequality $0\le1-\prod a_k\le\sum(1-a_k)$ gives
\[
0\le1-g(x)\le\sum_{k=1}^{\infty}\E(1-R_k)
=(1-x)\sum_{k=1}^{\infty}2^{-k}=1-x.
\]
Hence $g(x)\to1$ as $x\uparrow1$; the solution identically zero on $[0,1)$ is excluded.

Conditioning on the first draw, with bounded convergence for the subsequent infinite product, gives
\[
(1-x)g(x)=\int_x^1t g(t)\dd t.
\]
The integral of a bounded function first shows that $g$ is continuous on $[0,1)$, after which differentiation is valid. We obtain
\[
-g(x)+(1-x)g'(x)=-xg(x),\qquad g'(x)=g(x).
\]
Thus $g(x)=Ce^x$. The established condition $g(1-)=1$ implies $C=e^{-1}$, and therefore
\[
\boxed{\lim_{n\to\infty}\E\prod_{k=1}^nR_k=e^{-1}.}
\]

\answer{OP35}
Conditional on $H=h$, the remaining $n-1$ observations, before sorting, behave as independent uniform draws on $(0,h)$. Thus, for $0\le t\le h$,
\[
\Prob(L>t\mid H=h)=\left(1-\frac th\right)^{n-1}.
\]
Integrating the conditional survival function yields
\[
\E[L\mid H=h]=\int_0^h\left(1-\frac th\right)^{n-1}\dd t
=\frac hn.
\]
It follows that $\E L=\E H/n$ and
\[
\E[LH]=\E\!\left[H\E(L\mid H)\right]=\frac1n\E H^2,
\qquad \Cov(L,H)=\frac1n\Var(H).
\]
Replacing every observation by $1-X_i$ preserves the joint distribution and changes the minimum to $1-H$. Hence $L\overset d=1-H$, so $\Var(L)=\Var(H)>0$. Consequently,
\[
\boxed{\Corr(L,H)=\frac1n.}
\]
Conditioning on $H=h$ uses the version determined by the continuous densities, not a ratio of probabilities for a null event.

\answer{OP36}
Let the piece lengths be $D_1,\ldots,D_n$, with total length one. Positive lengths form a nondegenerate polygon if and only if the largest is strictly smaller than the sum of the others. Thus the event of interest is
\[
\max_iD_i<\frac12.
\]
Necessity follows from the strict triangle inequality along the other sides. For sufficiency, a convex construction also rules out self-intersections. Write $L$ for the largest length and $S$ for the total, with $2L<S$. On a circle of radius $r\ge L/2$, the smaller central angle for side $D_i$ is
\[
\theta_i(r)=2\arcsin\frac{D_i}{2r}.
\]
If the angle sum at $r=L/2$ is at least $2\pi$, continuity and its decrease to zero give a radius where the sum is exactly $2\pi$. Otherwise use the major arc $2\pi-\theta_L(r)$ for the longest side. We then need the other angles to sum to $\theta_L(r)$. Their sum minus $\theta_L(r)$ is negative at $r=L/2$, but for large $r$ it is asymptotic to $(S-2L)/r>0$. Continuity again gives the required radius. In either case, take these positive arcs successively around the circle and join neighboring vertices. This gives a convex polygon with the specified side lengths and no self-intersections.

The ordered uniform cut points have constant joint density $(n-1)!$. Transforming to spacings gives a constant density on the simplex $D_i>0$, $\sum D_i=1$, so the piece lengths are exchangeable. The first piece exceeds $1/2$ exactly when all cut points lie to the right of $1/2$. Hence
\[
\Prob(D_i>1/2)=\Prob(D_1>1/2)=2^{-(n-1)}.
\]
Two pieces cannot both exceed $1/2$, and equality has probability zero. Therefore,
\[
\boxed{\Prob(\text{a polygon can be formed})=1-\frac{n}{2^{n-1}}.}
\]
This gives $1/4$ for $n=3$ and $1/2$ for $n=4$.

\answer{OP37}
Sort the points as $X_{(1)}<\cdots<X_{(N)}$ and define the interior gaps
\[
D_i=X_{(i+1)}-X_{(i)},\qquad 1\le i\le N-1.
\]
Mutual nearest neighbors must be adjacent. Let $I_i$ indicate that the endpoints of adjacent edge $i$ choose one another. The first edge qualifies precisely when $D_1<D_2$, and the last edge has the analogous condition. An interior edge, $2\le i\le N-2$, qualifies precisely when
\[
D_i<D_{i-1},\qquad D_i<D_{i+1}.
\]
All uniform spacings, including the two boundary spacings, are exchangeable and have zero probability of ties. Each end edge therefore qualifies with probability $1/2$, while an interior gap is smallest among the relevant three with probability $1/3$.

Each point has a unique nearest neighbor, so distinct mutual pairs cannot share an endpoint. Every pair contributes exactly two good points. Thus, if $G$ is the number of good points,
\[
G=2\sum_{i=1}^{N-1}I_i,
\quad
\E G=2\left(\frac12+\frac12+\frac{N-3}{3}\right)
=\boxed{\frac{2N}{3}}.
\]
For $N=3$ the interior sum is empty and the same formula applies. The expected number of pairs, rather than points, is $N/3$.

